% SGA 3, Expose IV, English reconstruction.
% Scope: Propositions 4.4.9--4.4.10 and their proofs, stopping before
% Proposition 4.4.11.
% Authority: Polo--Gille Exp4-13oct24.pdf, component-local page 42
% (combined-reader page 230; printed pages 218 and 211).
% The external English Markdown was used only as a comparison.

\setcounter{footnote}{38}

\medskip
\noindent\SGARefTarget{sga3:iv:proposition:4:4:9}\textbf{命题 4.4.9.}
\textit{\(\widetilde{\mathscr C}\) 中每个等价关系都是普遍有效的
(\SGARefLink{sga3:iv:definition:3:3:3}{3.3.3}). 更确切地说,
设 \(R\) 是层 \(X\) 上的一个 \(\widetilde{\mathscr C}\)-等价关系.
则分离预层}
\[
  i(X)/i(R):S\longmapsto X(S)/R(S)
\]
\textit{的相伴层是 \(X\) 关于 \(R\) 的普遍有效商.}

用 \(X/R\) 表示 \(X\) 关于 \(R\) 的层商,
由 \SGARefLink{sga3:iv:theorem:4:4:1:item:ii:01}{4.4.1(ii)}, 它存在:
\[
  X/R=a\bigl(i(X)/i(R)\bigr).
\]
须证 \(X\to X/R\) 是普遍有效满态射, 且态射
\[
  f:R\longrightarrow X\times_{X/R}X
\]
是同构. 第一个断言已在 \SGARefLink{sga3:iv:proposition:4:4:3}{4.4.3} 中证明.
至于 \(f\), 它由对态射
\[
  i(R)\longrightarrow i(X)\times_{i(X/R)}i(X),
\]
施以函子 \(a\) 得到; 或者, 由于由
\SGARefLink{sga3:iv:lemma:4:4:5:item:ii:01}{4.4.5(ii)} 可知 \(i(X)/i(R)\) 分离,
从而
\[
  i(X)/i(R)\longrightarrow i(X/R)
\]
是单态射, 它由对典范态射
\[
  i(R)\longrightarrow
  i(X)\times_{\,i(X)/i(R)}i(X).
\]
施以该函子得到. 因此归结为在预层范畴中证明同一陈述.
而 \(i(X)/i(R)\) 是预层
\[
  S\longmapsto X(S)/R(S),
\]
故\SGARefLink{sga3:iv:proposition:4:4:9}{此断言}归结为其集合论的对应陈述,
后者是直接的.

\medskip
\noindent\SGARefTarget{sga3:iv:proposition:4:4:10}\textbf{命题 4.4.10.}
\textit{在 \SGARefLink{sga3:iv:proposition:4:4:9}{4.4.9} 的假设下,
设 \(Y\) 是 \(X\) 的子层, 用 \(R_Y\) 表示 \(R\)
在 \(Y\) 上诱导的等价关系. 则典范态射
(\SGARefLink{sga3:iv:definition:3:1:6}{3.1.6})}
\[
  Y/R_Y\longrightarrow X/R
\]
\textit{是单态射. 它将 \(Y/R_Y\) 与 \(X/R\) 的一个子层等同,
即复合}
\[
  Y\longrightarrow X\longrightarrow X/R.
\]
\textit{的层像.}

预层态射
\[
  i(Y)/i(R_Y)
  =i(Y)/i(R)_{i(Y)}
  \longrightarrow i(X)/i(R)
\]
是单态射. 由于函子 \(a\) 左正合
(\SGARefLink{sga3:iv:remark:4:3:16}{4.3.16}), 它将单态射映成单态射; 因而
\[
  Y/R_Y\longrightarrow X/R
\]
是单态射. 最后一个断言由交换图
\[
\begin{array}{ccc}
  Y&\longrightarrow&X\\
  \big\downarrow&&\big\downarrow\\
  Y/R_Y&\longrightarrow&X/R
\end{array}
\]
及 \(Y\to Y/R_Y\) 是覆盖态射这一事实得到.

鉴于\SGARefLink{sga3:iv:proposition:4:4:10}{此命题},
今后总将 \(Y/R_Y\) 与 \(X/R\) 的一个子层等同.

% Translation coverage: entire assigned source component.
% Source lines: 90; translated body lines before this audit comment: 92.
% Math preserved: 25 inline formulas, 12 display blocks.
% Reference targets: 2/2; links: 9/9; labels: 0/0; footnotes: 0/0.
% Unresolved TODOs: 0; compile and visual review delegated to the volume integrator.
